% ROZDZIAŁ 11 — do włączenia w main.tex
\chapter{Całki zależne od parametru}
\label{ch:parametry-calki}

\section{Problem i jego dokładne sformułowanie}

Niech $(X,\mathcal F,\mu)$ będzie przestrzenią z miarą, $T$ przestrzenią
parametrów, a $f:T\times X\to\mathbb C$ rodziną funkcji. Zajmujemy się
funkcją $F:T\to\mathbb C$ określoną wzorem
$F(t)=\int_X f(t,x)\,\mathrm d\mu(x)$, o ile $f(t,\cdot)$ jest całkowalna.
Pytamy, kiedy własności funkcji $t\mapsto f(t,x)$ przechodzą na $F$:
ciągłość, różniczkowalność, wyższa regularność i analityczność. Problem
polega na tym, że własność zachodząca osobno dla każdego $x$ nie
kontroluje automatycznie masy rozłożonej na całym $X$.

\begin{przyklad}[Koncentracja masy przy zerze]
\label{ex:params-escape-11}
Na $X=(0,\infty)$ z miarą Lebesgue'a połóżmy dla $t>0$
$f(t,x)=t\chi_{(0,1/t)}(x)$, a $f(0,x)=0$.
Wówczas $f(t,x)\to0$ dla każdego $x>0$ przy $t\downarrow0$,
ale $\int_X f(t,x)\,\mathrm d\lambda_1(x)=1$ dla każdego $t>0$.
Zatem $F$ nie jest ciągła w zerze, mimo ciągłości $t\mapsto f(t,x)$
w punkcie $0$ dla każdego $x$. Nie istnieje też jedna funkcja
$g\in L^1(X)$ dominująca wszystkie $f(t,\cdot)$ dla dostatecznie
małych $t>0$: w przeciwnym razie twierdzenie o zbieżności
zmajoryzowanej dawałoby sprzeczność.
\end{przyklad}

\begin{uwaga}[Zakres rozdziału]
W rozdziale 8 udowodniono podstawowe twierdzenia o całkach z parametrem.
Tutaj formułujemy ich wersje \emph{lokalne} i \emph{jednostajne},
wyprowadzamy twierdzenia wyższego rzędu i badamy kilka użytecznych
transformacji. Nie używamy twierdzenia Fubiniego, które zostanie
udowodnione dopiero w rozdziale~\ref{ch:produkty-miary}.
\end{uwaga}

\section{Dominacja lokalna i ciągłość w parametrze}

\begin{definicja}[Lokalna dominacja całkowalna]
\label{def:local-majorant-11}
Niech $(X,\mathcal F,\mu)$ będzie przestrzenią z miarą,
$(T,d)$ przestrzenią metryczną, $t_0\in T$ oraz
$f:T\times X\to\mathbb C$. Mówimy, że rodzina $f$ jest
\emph{lokalnie zdominowana całkowalnie w $t_0$}, jeżeli istnieją
otoczenie $U$ punktu $t_0$ w $T$ oraz mierzalna funkcja
$g:X\to[0,\infty)$ z $\int_Xg\,\mathrm d\mu(x)<\infty$ takie, że
istnieje $N\in\mathcal F$, $\mu(N)=0$, taki, że
$|f(t,x)|\le g(x)$ dla wszystkich $t\in U$ i $x\in X\setminus N$.
Warunek ten odróżniamy od słabszego wymagania, by dla każdego $t$
istniał osobny zbiór zerowy $N_t$. W twierdzeniu o ciągłości
słabszy wariant również wystarcza: dla ustalonego ciągu
$t_n\to t_0$ zbiór $N_{t_0}\cup\bigcup_n N_{t_n}$
pozostaje zerowy, toteż nadal można użyć twierdzenia
o zbieżności zmajoryzowanej. W pozostałych twierdzeniach
sprawdzamy potrzebny wariant oddzielnie.
\end{definicja}

\begin{twierdzenie}[Lokalne kryterium ciągłości]
\label{thm:local-continuity-11}
Niech $(X,\mathcal F,\mu)$ będzie przestrzenią z miarą,
$(T,d)$ przestrzenią metryczną, $t_0\in T$,
a $f:T\times X\to\mathbb C$ spełnia następujące warunki:
\begin{enumerate}[label=\textup{(\roman*)}]
\item dla każdego $t\in T$ funkcja $x\mapsto f(t,x)$ jest mierzalna;
\item dla $\mu$-prawie każdego $x$ odwzorowanie $t\mapsto f(t,x)$
      jest ciągłe w $t_0$;
\item istnieje całkowalna dominacja w pewnym otoczeniu $t_0$
      w sensie definicji~\ref{def:local-majorant-11}.
\end{enumerate}
Wtedy $F(t)=\int_X f(t,x)\,\mathrm d\mu(x)$ jest określona
w pewnym otoczeniu $t_0$ i ciągła w $t_0$.
Jeśli warunki te zachodzą w każdym $t_0\in T$, to $F$ jest
ciągła na $T$.
\end{twierdzenie}
\begin{proof}
Ustalmy otoczenie $U$ i funkcję $g$ z (iii). Dla $t\in U$
mamy $\int_X|f(t,x)|\,\mathrm d\mu(x)\le\int_Xg\,\mathrm d\mu(x)$,
więc $F(t)$ istnieje jako całka bezwzględnie zbieżna. Niech
$t_n\to t_0$. Dla dużych $n$ mamy $t_n\in U$, a funkcje
$f(t_n,\cdot)$ zbiegają do $f(t_0,\cdot)$ prawie wszędzie i
są zdominowane przez $g$. Twierdzenie o zbieżności zmajoryzowanej,
zastosowane do części rzeczywistej i urojonej lub bezpośrednio
do modułu różnicy, daje
$\int_X |f(t_n,x)-f(t_0,x)|\,\mathrm d\mu(x)\to0$.
W konsekwencji $|F(t_n)-F(t_0)|\to0$. Ciągowa charakteryzacja
ciągłości w przestrzeni metrycznej kończy dowód.
\end{proof}

\begin{wniosek}[Ciągłość w normie $L^1$]
\label{cor:l1-param-11}
Niech $(X,\mathcal F,\mu)$ będzie przestrzenią z miarą,
$(T,d)$ przestrzenią metryczną, $t_0\in T$,
a $f:T\times X\to\mathbb C$ rodziną o mierzalnych
przekrojach. Załóżmy, że $t\mapsto f(t,x)$ jest ciągła
w $t_0$ dla $\mu$-prawie każdego $x$ oraz że istnieją
otoczenie $U$ punktu $t_0$, mierzalny zbiór zerowy
$N\in\mathcal F$ i nieujemna całkowalna funkcja
$g:X\to[0,\infty)$ takie, że
$|f(t,x)|\le g(x)$ dla wszystkich $t\in U$
i $x\in X\setminus N$. Wtedy
\[
\int_X|f(t,x)-f(t_0,x)|\,\mathrm d\mu(x)\longrightarrow0
\qquad(t\to t_0).
\]
W szczególności odwzorowanie
$U\ni t\mapsto[f(t,\cdot)]\in L^1(\mu)$ jest
ciągłe w $t_0$.
\end{wniosek>

\begin{przyklad}[Brak potrzeby globalnej dominacji]
Niech $f(t,x)=e^{-tx}$ dla $x>0$ i $t>0$. Jeśli $t_0>0$,
to dla $t\in(t_0/2,3t_0/2)$ zachodzi
$0\le f(t,x)\le e^{-t_0x/2}$, a funkcja po prawej stronie
jest całkowalna na $(0,\infty)$. Zatem
$F(t)=\int_{(0,\infty)}e^{-tx}\,\mathrm d\lambda_1(x)=1/t$
jest ciągła na $(0,\infty)$. Nie istnieje jednak całkowalna
dominanta dla całej rodziny $\{f(t,\cdot):t>0\}$, gdyż
$\sup_{t>0} e^{-tx}=1$ dla każdego $x>0$.
\end{przyklad}

\section{Jednostajność zbieżności całek w parametrze}

\begin{twierdzenie}[Zbieżność jednostajna po scałkowaniu]
\label{thm:uniform-parameter-11}
Niech $(X,\mathcal F,\mu)$ będzie przestrzenią z miarą,
$T$ niepustym zbiorem, a $f_n,f:T\times X\to\mathbb C$
rodzinami funkcji o mierzalnych przekrojach.
Załóżmy, że istnieją mierzalne funkcje
$h_n,g:X\to[0,\infty)$, przy czym $g$ jest całkowalna, takie, że
$|f_n(t,x)-f(t,x)|\le h_n(x)\le g(x)$
dla wszystkich $t\in T$, wszystkich $n$ oraz prawie wszystkich
$x$ poza wspólnym zbiorem zerowym, a także $h_n(x)\to0$
prawie wszędzie. Załóżmy dodatkowo, że dla pewnego
$t_*\in T$ funkcja $f(t_*,\cdot)$ jest całkowalna oraz że
istnieje mierzalna funkcja całkowalna
$g_0:X\to[0,\infty)$ taka, że
$|f(t,x)-f(t_*,x)|\le g_0(x)$ dla każdego $t\in T$
i dla $\mu$-prawie każdego $x$ (zbiór wyjątkowy
może zależeć od $t$). Wtedy
\[
\sup_{t\in T}\left|\int_X f_n(t,x)\,\mathrm d\mu(x)
-\int_Xf(t,x)\,\mathrm d\mu(x)\right|\longrightarrow0.
\]
\end{twierdzenie}
\begin{proof}
Dla każdego $t\in T$ z nierówności dla całki wynika
$|\int_X(f_n-f)(t,x)\,\mathrm d\mu(x)|\le\int_Xh_n\,\mathrm d\mu(x)$.
Oszacowanie po prawej stronie nie zależy od $t$.
Twierdzenie Lebesgue'a o zbieżności zmajoryzowanej daje
$\int_Xh_n\,\mathrm d\mu(x)\to0$. Przechodząc do supremum
po $t$, otrzymujemy tezę.
\end{proof}

\begin{uwaga}[Mierzalność supremum nie jest automatyczna]
Jeśli chciałoby się wybrać $h_n(x)=\sup_{t\in T}|f_n(t,x)-f(t,x)|$,
trzeba najpierw sprawdzić mierzalność tej funkcji. Wystarcza
na przykład, aby $T$ było separowalną przestrzenią metryczną,
a odwzorowania $t\mapsto f_n(t,x)-f(t,x)$ były ciągłe dla
wszystkich $x$: wtedy supremum po $T$ równa się supremum
po ustalonym przeliczalnym podzbiorze gęstym. Jeśli nie
mamy tego warunku, lepiej założyć istnienie mierzalnej
majoranty $h_n$ bez odwoływania się do supremum.
\end{uwaga}

\section{Różniczkowanie pod znakiem całki}

\begin{twierdzenie}[Pierwsza pochodna — wersja lokalna]
\label{thm:local-differentiation-11}
Niech $(X,\mathcal F,\mu)$ będzie przestrzenią z miarą,
$I\subseteq\mathbb R$ przedziałem otwartym,
$t_0\in I$, a $f:I\times X\to\mathbb C$. Załóżmy, że:
\begin{enumerate}[label=\textup{(\roman*)}]
\item dla każdego $t\in I$ przekrój $f(t,\cdot)$ jest mierzalny;
\item dla prawie wszystkich $x$ funkcja $t\mapsto f(t,x)$ jest
      różniczkowalna na $I$;
\item istnieje $r>0$ z $[t_0-r,t_0+r]\subset I$ i funkcja
      $g\in L^1(\mu)$, $g\ge0$, taka, że
      $|\partial_t f(t,x)|\le g(x)$ dla wszystkich
      $t\in[t_0-r,t_0+r]$ poza wspólnym zbiorem zerowym;
\item funkcja $f(t_0,\cdot)$ jest całkowalna.
\end{enumerate}
Wówczas w otoczeniu $t_0$ funkcja
$F(t)=\int_Xf(t,x)\,\mathrm d\mu(x)$ jest dobrze określona,
różniczkowalna w $t_0$ i spełnia
\begin{equation}\label{eq:diff-param-11}
F'(t_0)=\int_X\partial_t f(t_0,x)\,\mathrm d\mu(x).
\end{equation}
Na wspólnym mierzalnym zbiorze wyjątkowym, na którym
nie zakładamy różniczkowalności, przyjmujemy
$\partial_t f(t_0,x)=0$.
\end{twierdzenie}
\begin{proof}
Poza wspólnym zbiorem zerowym ustalmy $x$ oraz $t$
i zastosujmy twierdzenie o wartości średniej do
rzeczywistej funkcji
$s\mapsto\operatorname{Re}(\overline{u}f(s,x))$,
gdzie $u\in\mathbb C$, $|u|=1$. Jej pochodna
ma moduł co najwyżej $g(x)$, zatem
$|\operatorname{Re}(\overline{u}(f(t,x)-f(t_0,x)))|
\le |t-t_0|g(x)$. Biorąc supremum po
$|u|=1$, otrzymujemy dokładnie
$|f(t,x)-f(t_0,x)|\le |t-t_0|g(x)$.
Argument nie wymaga ciągłości pochodnej ani
nieuzasadnionego zastosowania wzoru Newtona--Leibniza.
Stąd $f(t,\cdot)$ jest całkowalna dla $|t-t_0|<r$.
Dla $0<|h|<r$ iloraz
$q_h(x)=(f(t_0+h,x)-f(t_0,x))/h$ jest mierzalny
jako różnica funkcji mierzalnych; poza zbiorem zerowym
$q_h(x)\to\partial_tf(t_0,x)$ przy $h\to0$,
a $|q_h(x)|\le g(x)$. Pochodna $x\mapsto\partial_tf(t_0,x)$
jest mierzalna jako granica ilorazów różnicowych, na przykład
dla ciągu $h=1/n$ dostatecznie dużych $n$.
Twierdzenie o zbieżności zmajoryzowanej daje
$\int_Xq_h(x)\,\mathrm d\mu(x)\to
\int_X\partial_tf(t_0,x)\,\mathrm d\mu(x)$,
gdy $h\to0$. Z liniowości całki
$\int_Xq_h\,\mathrm d\mu(x)=(F(t_0+h)-F(t_0))/h$,
co jest dokładnie definicją pochodnej $F$.
\end{proof}

\begin{uwaga}[Dlaczego samo istnienie pochodnej nie wystarcza]
Dla $f(t,x)=t\chi_{(0,1/t)}(x)$ z pierwszego przykładu
nie można nawet otrzymać ciągłości całki w $t=0$.
Dominacja \emph{ilorazów różnicowych} lub lokalna dominacja
pochodnej są niezależnymi, istotnymi warunkami.
\end{uwaga}

\begin{twierdzenie}[Pochodne wyższego rzędu]
\label{thm:higher-differentiation-11}
Niech $(X,\mathcal F,\mu)$ będzie przestrzenią z miarą,
$I\subseteq\mathbb R$ otwartym przedziałem,
$m\ge1$ liczbą naturalną, a $f:I\times X\to\mathbb C$.
Załóżmy, że dla każdego $t\in I$ funkcja $f(t,\cdot)$ jest mierzalna,
a dla prawie wszystkich $x$ odwzorowanie $t\mapsto f(t,x)$
ma pochodne do rzędu $m$ włącznie, wszystkie ciągłe na $I$.
Załóżmy ponadto, że dla każdego zwartego przedziału
$J\subset I$ oraz każdego $1\le k\le m$
istnieje nieujemna mierzalna funkcja całkowalna
$g_{J,k}:X\to[0,\infty)$ taka, że
$|\partial_t^k f(t,x)|\le g_{J,k}(x)$ dla $t\in J$
poza wspólnym zbiorem zerowym, a $f(t_*,\cdot)$ jest
całkowalna dla pewnego $t_*\in I$.
Wtedy $F(t)=\int_Xf(t,x)\,\mathrm d\mu(x)$ należy
do $C^m(I)$ oraz dla wszystkich $t\in I$ i $1\le k\le m$
\begin{equation}\label{eq:higher-param-11}
F^{(k)}(t)=\int_X\partial_t^k f(t,x)\,\mathrm d\mu(x).
\end{equation}
Wszystkie pochodne podcałkowe określamy jako zero na
wspólnym mierzalnym zbiorze wyjątkowym, na którym
nie założono istnienia pochodnych.
\end{twierdzenie}
\begin{proof}
Dla $t\in I$ odcinek o końcach $t_*$ i $t$ jest zwarty
i zawiera się w $I$. Oszacowanie pierwszej pochodnej
zapewnia $|f(t,x)|\le|f(t_*,x)|+|t-t_*|g_{J,1}(x)$,
więc $f(t,\cdot)$ jest całkowalna. Dla ustalonego
$t\in I$ zastosujmy twierdzenie~\ref{thm:local-differentiation-11}
w małym otoczeniu $t$: otrzymujemy istnienie $F'$ i wzór
z $k=1$. Gdy $1\le k<m$, funkcja $\partial_t^kf(t,\cdot)$
jest mierzalna jako punktowa granica ilorazów różnicowych.
Dzięki majorantom $g_{J,k}$ jest całkowalna. Stosując
twierdzenie~\ref{thm:local-differentiation-11} do
$\partial_t^kf$, otrzymujemy pochodną $(k+1)$-szego rzędu.
Ciągłość każdej funkcji $t\mapsto\int_X\partial_t^kf(t,x)\,\mathrm d\mu(x)$
wynika z twierdzenia~\ref{thm:local-continuity-11},
gdyż $\partial_t^kf(t,x)$ jest punktowo ciągła w $t$
i lokalnie zdominowana. Zatem $F\in C^m(I)$.
\end{proof}

\begin{przyklad}[Całka Laplace'a i pochodne]
\label{ex:laplace-11}
Dla $t>0$ kładziemy $F(t)=\int_{(0,\infty)}e^{-tx}\,\mathrm d\lambda_1(x)$.
Mamy $F(t)=1/t$. Ponieważ
$\partial_t^k e^{-tx}=(-x)^k e^{-tx}$, dla
$t$ w zwartym podprzedziale $(0,\infty)$ funkcja
$x^k e^{-tx}$ jest zdominowana przez
$x^ke^{-ax}$, gdzie $a>0$ zależy od podprzedziału.
Ta ostatnia funkcja jest całkowalna, więc
\[
\int_{(0,\infty)}x^ke^{-tx}\,\mathrm d\lambda_1(x)
=\frac{k!}{t^{k+1}}\qquad(k\in\mathbb N_0).
\]
W szczególności wzór na momenty rozkładu wykładniczego
wynika z różniczkowania pod znakiem całki, a nie z
wielokrotnego całkowania przez części.
\end{przyklad}

\section{Zbieżność według miary zamiast dominacji punktowej}

\begin{twierdzenie}[Kryterium Vitaliego dla parametru]
\label{thm:param-vitali-11}
Niech $(X,\mathcal F,\mu)$ będzie przestrzenią z miarą
skończoną, $(T,d)$ przestrzenią metryczną, $t_0\in T$,
a $f:T\times X\to\mathbb C$ rodziną funkcji o mierzalnych
przekrojach. Załóżmy, że istnieje otoczenie $U$ punktu
$t_0$ takie, że rodzina modułów
$\{|f(t,\cdot)|:t\in U\}$ jest jednostajnie całkowalna
w sensie Definicji~\ref{def:uniform-integrability}
oraz dla każdego $\varepsilon>0$ zachodzi
\[
\mu(\{x\in X:|f(t,x)-f(t_0,x)|>\varepsilon\})
\longrightarrow0\qquad(t\to t_0).
\]
Jest to naturalne rozszerzenie zbieżności według miary
z Definicji~\ref{def:convergence-measure} na funkcje
zespolone. Wówczas
$\|f(t,\cdot)-f(t_0,\cdot)\|_{L^1(\mu)}\to0$,
a więc funkcja $t\mapsto\int_X f(t,x)\,\mathrm d\mu(x)$
jest ciągła w $t_0$.
\end{twierdzenie}
\begin{proof}
Załóżmy przeciwnie, że granica norm nie jest zerem.
Istnieją wówczas $\varepsilon>0$ i ciąg $t_n\to t_0$,
$t_n\in U$, dla którego
$\|f(t_n,\cdot)-f(t_0,\cdot)\|_1\ge\varepsilon$.
Funkcje $f(t_n,\cdot)$ zbiegają według miary do
$f(t_0,\cdot)$, a rodziny ich części rzeczywistych i urojonych
są jednostajnie całkowalne. Twierdzenie Vitaliego
z rozdziału~\ref{ch:twierdzenia-graniczne}
(dowodzone osobno dla części rzeczywistych i urojonych)
implikuje zbieżność w $L^1$, sprzeczną z wybraną
nierównością. W szczególności wszystkie całki
są określone i ich ciągłość wynika z nierówności
$|\int_X(f(t,x)-f(t_0,x))\,\mathrm d\mu(x)|
\le\|f(t,\cdot)-f(t_0,\cdot)\|_1$.
\end{proof}

\begin{uwaga}
Jednostajna całkowalność nie wymaga wspólnej
całkowalnej dominanty. W przeciwieństwie do twierdzenia
o dominacji lokalnej założenie o zbieżności jest tutaj
formułowane według miary, a nie punktowo. Dla $\mu(X)=\infty$
potrzebna jest dodatkowa kontrola ogonów; sam warunek
jednostajnej całkowalności na małych zbiorach nie wystarcza.
\end{uwaga}

\section{Zmienna dziedzina całkowania}

\begin{twierdzenie}[Całkowanie po zależnej od parametru dziedzinie]
\label{thm:moving-domains-11}
Niech $(X,\mathcal F,\mu)$ będzie przestrzenią z miarą,
$(T,d)$ przestrzenią metryczną, $t_0\in T$,
a $A_t\in\mathcal F$ dla każdego $t\in T$.
Niech $u:T\times X\to\mathbb C$
ma mierzalne przekroje $u(t,\cdot)$ i załóżmy, że dla prawie
wszystkich $x$ zachodzi
$\chi_{A_t}(x)u(t,x)\to\chi_{A_{t_0}}(x)u(t_0,x)$
dla $t\to t_0$. Jeżeli w pewnym otoczeniu $U$ punktu $t_0$
istnieje nieujemna funkcja $g\in L^1(\mu)$ taka, że
\emph{dla każdego} $t\in U$ nierówność
$|\chi_{A_t}(x)u(t,x)|\le g(x)$
zachodzi dla $\mu$-prawie każdego $x$ (zbiór wyjątkowy
może zależeć od $t$), to
$t\mapsto\int_{A_t}u(t,x)\,\mathrm d\mu(x)$
jest ciągła w $t_0$.
\end{twierdzenie}
\begin{proof}
Definiujemy na całym $X$ funkcję
$f(t,x)=u(t,x)$ dla $x\in A_t$ i $f(t,x)=0$
dla $x\notin A_t$. Jest ona mierzalna w $x$ dla każdego $t$;
wyrażenia $\chi_{A_t}u$ należy tak interpretować,
aby uniknąć nieoznaczonych iloczynów $0\cdot\infty$.
Ustalmy ciąg $t_n\to t_0$. Po pominięciu skończenie
wielu wyrazów możemy założyć $t_n\in U$. Dla każdego
$n\ge0$ wybierzmy zbiór zerowy $N_n$, poza którym
$|f(t_n,x)|\le g(x)$; w tym zapisie $t_0$ odpowiada
$n=0$. Zbiór $N=\bigcup_{n\ge0}N_n$ ma miarę zero.
Poza $N$ oraz zbiorem wyjątkowym z założenia o zbieżności
punktowej zachodzi jednocześnie $f(t_n,x)\to f(t_0,x)$
i $|f(t_n,x)|\le g(x)$ dla wszystkich $n\ge0$.
Twierdzenie Lebesgue'a o zbieżności zmajoryzowanej
zapewnia więc zbieżność całek wzdłuż każdego ciągu
$t_n\to t_0$. Ciągowa charakteryzacja ciągłości
w przestrzeni metrycznej kończy dowód. W ten sposób
nie użyliśmy silniejszego, niewymaganego założenia
o wspólnym zbiorze wyjątkowym dla całej rodziny parametrów.
\end{proof}

\begin{przyklad}[Przedziały ruchome i brak dominacji]
Niech $u(t,x)=1$ na $(0,\infty)$ i $A_t=(0,t)$ dla $t>0$.
Dla $t$ w otoczeniu $t_0>0$ wszystkie $A_t$ mieszczą się
w $(0,2t_0)$; funkcja $g=\chi_{(0,2t_0)}$ jest
całkowalną dominantą, a dla każdego $x\ne t_0$
zachodzi $\chi_{A_t}(x)\to\chi_{A_{t_0}}(x)$.
Stąd $\int_{A_t}1\,\mathrm d\lambda_1(x)=t$ zależy ciągle od $t$.
Ten sam argument przestaje działać, jeśli dziedziny rozrastają
się tak, że masa ucieka poza każdy zbiór skończonej miary.
\end{przyklad}

\section{Całki niewłaściwe i jednostajna kontrola ogonów}

\begin{twierdzenie}[Kryterium jednostajnej kontroli ogonów]
\label{thm:tails-11}
Niech $t_0\in T$, gdzie $(T,d)$ jest przestrzenią metryczną,
a $(X,\mathcal F,\mu)$ przestrzenią z miarą. Załóżmy, że
istnieje wyczerpanie $X_1\subseteq X_2\subseteq\dots$,
$\bigcup_{n\ge1}X_n=X$, zbiory $X_n\in\mathcal F$
o skończonej mierze, oraz $f:T\times X\to\mathbb C$,
dla której $f(t,\cdot)$ jest całkowalna na każdym $X_n$.
Załóżmy ponadto, że $t\mapsto\int_{X_n}f(t,x)\,\mathrm d\mu(x)$
jest ciągła w $t_0$ dla każdego $n$, oraz że istnieje
otoczenie $U$ punktu $t_0$ takie, że
\begin{equation}\label{eq:uniform-tails-11}
\lim_{n\to\infty}\sup_{t\in U}
\int_{X\setminus X_n}|f(t,x)|\,\mathrm d\mu(x)=0.
\end{equation}
Wtedy wszystkie $f(t,\cdot)$ dla $t\in U$ są całkowalne,
a $t\mapsto\int_Xf(t,x)\,\mathrm d\mu(x)$ jest ciągła w $t_0$.
\end{twierdzenie}
\begin{proof}
Warunek \eqref{eq:uniform-tails-11} gwarantuje, że dla
pewnego $n$ i wszystkich $t\in U$ całki modułów po
$X\setminus X_n$ są skończone; ich suma z całką po $X_n$
jest więc skończona. Ustalmy $\varepsilon>0$ i wybierzmy
$n$ tak duże, aby supremum w
\eqref{eq:uniform-tails-11} było mniejsze niż $\varepsilon/3$.
Dla $t\in U$ mamy
\begin{align*}
\left|\int_Xf(t,x)\,\mathrm d\mu(x)-\int_Xf(t_0,x)\,\mathrm d\mu(x)\right|
&\le\left|\int_{X_n}(f(t,x)-f(t_0,x))\,\mathrm d\mu(x)\right|\\
&\quad+\int_{X\setminus X_n}|f(t,x)|\,\mathrm d\mu(x)
+\int_{X\setminus X_n}|f(t_0,x)|\,\mathrm d\mu(x).
\end{align*}
Dwa ostatnie składniki są mniejsze niż $2\varepsilon/3$.
Pierwszy jest mniejszy niż $\varepsilon/3$, gdy $t$
jest dostatecznie blisko $t_0$ dzięki ciągłości
całki na $X_n$. To kończy dowód.
\end{proof}

\begin{uwaga}[Ograniczenie kryterium]
Warunek \eqref{eq:uniform-tails-11} jest wystarczający, lecz
niekonieczny dla ciągłości całki: przy funkcjach zmieniających
znak mogą występować znoszące się wkłady z ogona.
Kryterium jest użyteczne zwłaszcza dla całek niewłaściwych,
gdzie trudno znaleźć jedną majorantę całkowalną na całym $X$.
\end{uwaga}

\section{Transformata Laplace'a i funkcje holomorficzne}

\begin{twierdzenie}[Holomorficzność transformaty Laplace'a]
\label{thm:laplace-holo-11}
Niech $a\in\mathbb R$ i niech $h:(0,\infty)\to\mathbb C$
będzie mierzalna oraz
$\int_{(0,\infty)}e^{-\sigma x}|h(x)|\,\mathrm d\lambda_1(x)<\infty$
dla każdego $\sigma>a$. Dla $z\in\mathbb C$,
$\operatorname{Re}z>a$, określmy
$L_h(z)=\int_{(0,\infty)}e^{-zx}h(x)\,\mathrm d\lambda_1(x)$.
Wtedy $L_h$ jest holomorficzna w półpłaszczyźnie
$\{\operatorname{Re}z>a\}$ i dla każdego $k\in\mathbb N_0$
\begin{equation}\label{eq:laplace-derivative-11}
L_h^{(k)}(z)=\int_{(0,\infty)}(-x)^k e^{-zx}h(x)
\,\mathrm d\lambda_1(x).
\end{equation}
\end{twierdzenie}
\begin{proof}
Ustalmy $z_0$ z $\operatorname{Re}z_0>a$ i liczbę
$\sigma$ spełniającą $a<\sigma<\operatorname{Re}z_0$.
W pewnym otoczeniu $z_0$ zachodzi $\operatorname{Re}z>\sigma$.
Dla $k\ge0$ istnieje $C_{k,\sigma,z_0}<\infty$ taka, że
$x^k e^{-\operatorname{Re}z\,x}\le
C_{k,\sigma,z_0}e^{-\sigma x}$ w nieco mniejszym
otoczeniu $z_0$: wybieramy pośrednią liczbę
$\sigma'>\sigma$ z $\operatorname{Re}z>\sigma'$
i korzystamy z ograniczoności $x^ke^{-(\sigma'-\sigma)x}$.
Zatem całki po prawej stronie
\eqref{eq:laplace-derivative-11} są zbieżne bezwzględnie.

Aby udowodnić różniczkowalność zespoloną, dla małego
$w\in\mathbb C\setminus\{0\}$ rozważamy
\[
\frac{e^{-(z_0+w)x}-e^{-z_0x}}{w}
=-x\int_{[0,1]}e^{-(z_0+sw)x}\,\mathrm d\lambda_1(s).
\]
Tożsamość z prawą stroną jest klasyczną
całką po rzeczywistym odcinku $s\in[0,1]$,
a nie zamianą kolejności dwóch całkowań.
Moduł ilorazu jest ograniczony przez
$x e^{-\sigma'x}$ dla odpowiednio małych $w$,
co po pomnożeniu przez $|h(x)|$ jest zdominowane
przez stałą razy $e^{-\sigma x}|h(x)|$.
Dla każdego $x$ iloraz dąży do $-xe^{-z_0x}$.
Twierdzenie o zbieżności zmajoryzowanej daje
$L_h'(z_0)=-\int_{(0,\infty)}xe^{-z_0x}h(x)
\,\mathrm d\lambda_1(x)$. Powtarzając argument dla
$(-x)^kh(x)$, otrzymujemy wszystkie pochodne.
Ponieważ pochodna zespolona istnieje w każdym punkcie
półpłaszczyzny, $L_h$ jest holomorficzna.
\end{proof}

\begin{przyklad}[Funkcja Gamma]
\label{ex:gamma-param-11}
Dla $s>0$ definiujemy
$\Gamma(s)=\int_{(0,\infty)}x^{s-1}e^{-x}\,\mathrm d\lambda_1(x)$.
Całka jest skończona: przy zerze $x^{s-1}$ jest
całkowalne dla $s>0$, a w nieskończoności czynnik
$e^{-x}$ dominuje każdą potęgę $x$.
Dla każdego $s_0>0$ wybieramy $0<a<s_0<b$;
na $s\in[a,b]$ pochodne rzędu $k$ funkcji podcałkowej
wynoszą $x^{s-1}(\log x)^k e^{-x}$ i mają całkowalną
majorantę
\[
\chi_{(0,1]}(x)x^{a-1}|\log x|^k
+\chi_{(1,\infty)}(x)x^{b-1}(1+\log x)^k e^{-x}.
\]
W konsekwencji $\Gamma\in C^\infty((0,\infty))$ i
$\Gamma^{(k)}(s)=\int_{(0,\infty)}
 x^{s-1}(\log x)^k e^{-x}\,\mathrm d\lambda_1(x)$.
Całkując przez części na $[\varepsilon,R]$ i przechodząc
do granic $\varepsilon\downarrow0$, $R\uparrow\infty$,
dostajemy $\Gamma(s+1)=s\Gamma(s)$.
Ponadto $\Gamma(1)=1$, więc $\Gamma(n+1)=n!$ dla $n\ge0$.
\end{przyklad}

\section{Podsumowanie}
Dominacja lokalna pozwala przenosić ciągłość i
różniczkowalność pod znak całki bez narzucania
jednej globalnej majoranty. Jednostajna kontrola
błędów oraz ogonów pozwala badać rodziny funkcji,
dla których pojedyncza dominanta bywa niedostępna.
Wersja zespolona prowadzi do holomorficzności
transformaty Laplace'a i umożliwia obliczanie jej
pochodnych. W kolejnym rozdziale skonstruujemy
miary produktowe, dzięki którym można zamieniać
kolejność całkowania względem dwóch niezależnych
miar.
